Code général des impôts ; LFSS et LF 2026 ; calculs de l'auteur.\par Note : impôt total sur 100 de bénéfice avant impôt, droit de 2026 ; entre parenthèses, le taux du prélèvement personnel sur le dividende. Impôt sur les sociétés de 25 \% plus la contribution sociale de 3,3 \% (25,8 \%), puis prélèvement personnel sur le dividende : PFU de 31,4 \% (12,8 \% d'impôt sur le revenu et 18,6 \% de prélèvements sociaux), majoré de la contribution exceptionnelle sur les hauts revenus (4 \%) ou porté au plancher de 20 \% de la contribution différentielle. Holding : dividende exonéré à 95 \% (régime mère-fille), soit 1,25 \% d'impôt. Salaire : impôt sur le revenu à 45 \% et CEHR à 4 \%, plus CSG et CRDS (9,7 \%, dont 6,8 points déductibles), hors cotisations sociales. Entreprises de plus de 3 Md\euro{} de chiffre d'affaires (taux le plus élevé de la contribution exceptionnelle, dont les quelque 450 redevables comprennent aussi des groupes de 1 à 3 Md\euro{}) : 56 à 61 \% selon la situation de l'actionnaire.
